% 2 Background (~1.0 p). The analyzers are the subjects of the study, not its
% contribution, so they are kept apart from the technique.
\section{Background}
\label{sec:background}

\subsection{JPAMB and Outcome Predictions}
\label{sec:jpamb}
JPAMB (Java Program Analysis Micro Benchmarks)~\citep{kalhauge2026jpamb} is a suite
of small static Java methods, each annotated with test cases that state which
outcome an input produces, for example that input \texttt{(0)} leads to a
division by zero. An analyzer receives a method identifier and the bytecode,
and answers, for each of six outcomes, whether some input reaches it: normal
termination (\texttt{ok}), divide by zero, assertion error, out of bounds,
null pointer, and non-termination (\texttt{*}). We call the outcomes it
answers \emph{yes} to its \emph{outcome set} for the method. An analyzer is
sound for a method if the outcome set contains every outcome some input
reaches, and precise if it contains no other. JPAMB scores the answers against
the outcomes of the annotated cases.

\subsection{Metamorphic Testing}
\label{sec:mt}
Metamorphic testing replaces a missing oracle by a \emph{metamorphic
relation} between the outputs for related inputs~\citep{chen2018}. Our system
under test is an analyzer and its inputs are methods. The relation is: two
methods with the same behavior on every input get the same outcome set. A
sound and fully precise analyzer satisfies it trivially, because both outcome
sets equal the true one. A sound but imprecise analyzer over-approximates, and
how much it over-approximates may depend on how the method is written. A
violation, an inconsistency, shows that one of the two outcome sets is
imprecise or, if a reachable outcome is missing, unsound (\cref{sec:inconsistency} defines how we count them). The relation is a
robustness property rather than a correctness requirement.

\subsection{The Three Analyzers Under Test}
\label{sec:analyzers}
The three analyzers are abstract interpreters~\citep{cousot1977} over JPAMB
bytecode. They share one implementation and differ only in the abstract domain
for integers and in whether they refine values at conditional jumps
(\cref{tab:analyzers}). Each computes a fixpoint over the method's
control-flow graph with a worklist. The outcome set collects what the reachable
abstract states can produce: \texttt{ok} at a return; divide by zero at a
\texttt{/} or \texttt{\%} whose divisor may be 0; null pointer at an array
access whose reference may be null, and out of bounds at every array access
whose reference may be non-null, because no domain tracks array lengths;
assertion error at a \texttt{throw}; and \texttt{*} at every reachable
backward jump, so there is no termination analysis. The analysis is
intraprocedural: a call is not analyzed, may raise an assertion error, and
returns an unknown value. References and booleans share the integer domain
(null is 0).

\begin{table}[t]
\caption{The three analyzers under test.}
\label{tab:analyzers}
\small
\begin{tabular}{@{}llll@{}}
\toprule
Analyzer & Domain & Refinement & Widening \\
\midrule
\texttt{static-sign} & subsets of $\{-,0,+\}$ & no & join (finite) \\
\texttt{static-sign-narrow} & subsets of $\{-,0,+\}$ & yes & join (finite) \\
\texttt{static-interval} & intervals $[l,h]$ of 32-bit ints & yes & at loop heads \\
\bottomrule
\end{tabular}
\end{table}

\paragraph{Domains}
Both domains abstract sets of 32-bit integers ordered by inclusion, and both
form Galois connections with them: for signs,
$\alpha(V)=\{\mathrm{sign}(v) \mid v\in V\}$ and
$\gamma(S)=\{v \mid \mathrm{sign}(v)\in S\}$; for intervals,
$\alpha(V)=[\min V,\max V]$ and $\gamma([l,h])=\{v \mid l\le v\le h\}$. We
write $\pm\infty$ for a bound at the 32-bit limit. For \texttt{+}, \texttt{-},
\texttt{*} and \texttt{/}, interval arithmetic computes the hull of the
results at the bounds, and a result that leaves the 32-bit range, where Java
wraps around, becomes $[-\infty,\infty]$; \texttt{\%} uses $|a \% b| < |b|$,
and other operators give $[-\infty,\infty]$. The sign domain ignores
wrap-around, so its predictions are unsound when an operation overflows.

\paragraph{Refinement at branches}
At a conditional jump such as \texttt{if (x > 0)}, an analyzer with refinement
intersects the abstract value of each compared local variable with the values
that make the condition true on one edge and false on the other; a compared
expression such as \texttt{x+1} is not refined. With or without refinement, an
edge is dropped when the comparison cannot have that outcome, but without
refinement its successor receives the unrefined state. This is also called
narrowing at branches, hence the name \texttt{static-sign-narrow}; it is not
the narrowing operator that recovers precision after
widening~\citep{moller2018spa}, which none of the analyzers uses.

\paragraph{Widening}
Every cycle in the control-flow graph contains an edge to an equal or smaller
instruction offset, and the analyzer widens at its target. States at a node
only grow, and at a widening point every strict increase of a bound sets it to
$\pm\infty$, so each widening point changes finitely often and the iteration
terminates. For signs, widening is the join, which suffices because the
lattice is finite; for intervals, ascending chains have length about $2^{33}$,
so iteration without widening could take that many steps.

\paragraph{Neither domain is uniformly more precise}
An interval has no holes: refining $[-\infty,\infty]$ with \texttt{x != 0}
leaves $[-\infty,\infty]$, since \texttt{!=} removes a constant only at an
endpoint (e.g.\ $[0,h]$ becomes $[1,h]$), while \texttt{static-sign-narrow}
obtains $\{-,+\}$ and proves a later \texttt{1 / x} safe. Conversely,
\texttt{static-interval} proves \texttt{if (n > 100) \{ if (n < 50) ... \}}
infeasible, which signs cannot. Of the \NumMethods{} benchmark methods,
\texttt{static-sign-narrow} predicts a strictly smaller outcome set than
\texttt{static-interval} on \NumNarrowBeatsInterval{} and a strictly larger one
on \NumIntervalBeatsNarrow{}, so the analyzers differ in precision for reasons
unrelated to where a rewrite is placed.
